% =============================================================================
% htgr-prke.tex
% Point Reactor Kinetics Equations with HTGR Sim v1
%
% Adapted from the original FHR PRKE teaching presentation.
%
% Compile:
%   lualatex htgr-prke.tex
%   biber htgr-prke
%   lualatex htgr-prke.tex
%   lualatex htgr-prke.tex
% =============================================================================

\documentclass[aspectratio=169]{beamer}

\usepackage{beamerthememidcenturymodern}
\usepackage[
    backend=biber,
    style=authoryear,
    sorting=nyt
]{biblatex}

\addbibresource{demo.bib}

% -----------------------------------------------------------------------------
% Theme
% -----------------------------------------------------------------------------

\mcmTheme{Kraft}

\usepackage{graphicx}
\usepackage{amsmath}
\usepackage{booktabs}

\AtBeginSection[]{%
    \begin{frame}[plain]
        \sectionpage
    \end{frame}
}

\AtBeginSubsection[]{%
    \begin{frame}[plain]
        \subsectionpage
    \end{frame}
}

\AtEveryCite{\color{mcmPrimary}}

% -----------------------------------------------------------------------------
% Metadata
% -----------------------------------------------------------------------------

\institute{Singapore Nuclear Research and Safety Institute}

\title[PRKE with HTGR Simulator]
      {Point Reactor Kinetics Equations (PRKE) with HTGR Simulator}

\subtitle{An Interactive Introduction to Reactor Dynamics}

\author{Dr Theodore Ong, SNRSI, NUS}

\date{\today}

% -----------------------------------------------------------------------------
% Screenshot filenames
%
% Replace these with screenshots exported from HTGR Sim v1.
% If the files do not exist, the deck will display placeholder boxes.
% -----------------------------------------------------------------------------

\newcommand{\htgrfull}{htgr_sim_full.png}
\newcommand{\htgrcontrols}{htgr_sim_controls.png}
\newcommand{\htgrshutdown}{htgr_sim_shutdown.png}
\newcommand{\htgrshutdownplot}{htgr_sim_shutdown_plot.png}


\begin{document}


% =============================================================================
% TITLE
% =============================================================================

\begin{frame}
    \titlepage
\end{frame}


% =============================================================================
% PART I -- OPERATE FIRST
% =============================================================================

\section{Operate First}


% -----------------------------------------------------------------------------
% Mission
% -----------------------------------------------------------------------------

\begin{frame}{Mission Briefing: Change Reactor Power}

    \begin{columns}

        \begin{column}{0.58\textwidth}

            The HTR-10 demonstration plant is running steadily at
            \textbf{about 10\,MWth}, near its rated thermal power.

            \vspace{1em}

            \textbf{The Situation:}

            The grid requests a change in power.

            \vspace{1em}

            \begin{alertblock}{Your Objective}
                Raise reactor thermal power from \textbf{about 10\,MWth} to
                \textbf{20\,MWth} in a controlled manner, and stabilise the
                plant there.
            \end{alertblock}

            \vspace{0.5em}

            Watch not only reactor power, but also:

            \begin{itemize}
                \item reactivity,
                \item fuel/core temperature,
                \item helium temperatures and flow,
                \item steam-generator response.
            \end{itemize}

        \end{column}

        \begin{column}{0.42\textwidth}

            \centering

            \IfFileExists{\htgrfull}
            {
                \includegraphics[width=\textwidth]{\htgrfull}
            }
            {
                \fbox{
                    \parbox[c][0.48\textheight][c]{0.9\textwidth}{
                        \centering
                        Insert HTGR Sim v1 full-interface screenshot here.
                    }
                }
            }

        \end{column}

    \end{columns}

\end{frame}


% -----------------------------------------------------------------------------
% Constraints
% -----------------------------------------------------------------------------

\begin{frame}{Operational Constraints}

    A reactor is not a volume knob.

    A reactivity change alters the neutron population, while thermal feedback
    and delayed neutrons determine how the response develops.

    \vspace{1em}

    \begin{block}{Rules}

        \begin{itemize}

            \item Use the simulator's reactivity/control input to change power.

            \item Aim for \textbf{20\,MWth}, approaching it without a
                  large overshoot.

            \item Observe reactivity in dollars as well as thermal power.

            \item Do not treat this demonstrator as an operational reactor model.

        \end{itemize}

    \end{block}

    \begin{alertblock}{Demonstration Model}

        HTGR Sim v1 is an educational/research demonstrator.

        Plant coefficients, controller constants and some operating assumptions
        are illustrative.

    \end{alertblock}

\end{frame}


% -----------------------------------------------------------------------------
% Safety limit
% -----------------------------------------------------------------------------

\begin{frame}{Safety Limit}

    \begin{alertblock}{Failure Condition}

        If core thermal power exceeds \textbf{500\,MWth} at any point during
        the manoeuvre, you have lost control of the reactor and the exercise
        counts as failed.

    \end{alertblock}

    \vspace{1em}

    That is roughly fifty times rated power, so it is not a near miss.
    Reaching it means the excursion got away from you entirely.

    \vspace{1em}

    A real HTR-10 would never get there. Its protection system trips on high
    neutron flux and on high core outlet temperature, and the fuel temperature
    limit of 1230\,\textdegree C governs long before power does. This limit is a
    teaching marker for the exercise, not a plant safety limit.

\end{frame}


% -----------------------------------------------------------------------------
% Interface
% -----------------------------------------------------------------------------

\begin{frame}{Your Interface: HTGR Sim v1}

    \centering

    \IfFileExists{\htgrcontrols}
    {
        \includegraphics[
            width=0.92\textwidth,
            height=0.72\textheight,
            keepaspectratio
        ]{\htgrcontrols}
    }
    {
        \fbox{
            \parbox[c][0.58\textheight][c]{0.82\textwidth}{
                \centering
                Insert HTGR Sim v1 control-panel screenshot here.
            }
        }
    }

\end{frame}


% -----------------------------------------------------------------------------
% Challenge
% -----------------------------------------------------------------------------

\begin{frame}{The Challenge Begins}

    \centering

    \Huge
    \textbf{Volunteers Needed!}

    \vspace{1.2em}

    \normalsize

    Who has the steady hand to take the reactor from about 10\,MWth to
    \textbf{20\,MWth} and hold it there, without ever passing
    \textbf{500\,MWth}?

    \vspace{1.2em}

    \hrule

    \vspace{0.8em}

    \textbf{Observers:}

    \begin{itemize}

        \item What happens immediately after positive reactivity is inserted?

        \item Does power stop changing immediately when the operator stops
              moving the control?

        \item How quickly does temperature respond compared with neutron power?

        \item Does increasing temperature help or oppose the power rise?

    \end{itemize}

\end{frame}


% -----------------------------------------------------------------------------
% Debrief
% -----------------------------------------------------------------------------

\begin{frame}{Debrief: Intuition vs.\ Reactor Dynamics}

    \begin{enumerate}

        \item Why did the power respond much faster than the plant temperatures?

        \item Why was it easy to overshoot?

        \item Why does the response have memory even after the control input
              changes?

        \item Why does increasing temperature tend to oppose the excursion
              in this model?

    \end{enumerate}

    \vspace{1em}

    \begin{center}

        \textbf{
            To explain what we observed, we need to follow the neutron
            population.
        }

    \end{center}

\end{frame}


% =============================================================================
% PART II -- PRKE
% =============================================================================

\section{From Neutron Balance to PRKE}


% -----------------------------------------------------------------------------
% Neutron balance
% -----------------------------------------------------------------------------

\begin{frame}{Step 1: The Neutron Balance}

    At its simplest, reactor power follows the neutron population $n$
    \cite{bell1970nuclear}:

    \[
        \text{Rate of change}
        =
        \text{Production}
        -
        \text{Loss}
    \]

    \[
        \frac{dn}{dt}
        =
        \text{Production}
        -
        \text{Loss}
    \]

    \vspace{1em}

    The detailed balance contains:

    \begin{itemize}

        \item fission,

        \item absorption,

        \item leakage,

        \item moderation,

        \item energy-dependent neutron interactions.

    \end{itemize}

    Point kinetics compresses the spatial reactor into one time-dependent
    amplitude.

\end{frame}


% -----------------------------------------------------------------------------
% k-effective
% -----------------------------------------------------------------------------

\begin{frame}{What Does $k_{\mathrm{eff}}$ Tell Us?}

    The effective multiplication factor describes the change in neutron
    population from one generation to the next:

    \[
        k_{\mathrm{eff}}
        =
        \frac{
            \text{neutrons in generation }(n+1)
        }{
            \text{neutrons in generation }n
        }
    \]

    \vspace{1em}

    \begin{itemize}

        \item $k_{\mathrm{eff}} < 1$:
              \textbf{subcritical}

        \item $k_{\mathrm{eff}} = 1$:
              \textbf{critical}

        \item $k_{\mathrm{eff}} > 1$:
              \textbf{supercritical}

    \end{itemize}

    \vspace{1em}

    We normally express departure from criticality using \textbf{reactivity}:

    \[
        \rho
        =
        \frac{k_{\mathrm{eff}}-1}{k_{\mathrm{eff}}}
    \]

\end{frame}


% -----------------------------------------------------------------------------
% Prompt kinetics
% -----------------------------------------------------------------------------

\begin{frame}{Step 2: Prompt-Neutron Kinetics}

    If every fission neutron appeared promptly, the neutron population would
    respond on the prompt-generation timescale.

    \vspace{1em}

    A useful schematic form is:

    \[
        \frac{dn}{dt}
        \sim
        \frac{\rho}{\Lambda} n
    \]

    where $\Lambda$ is the prompt neutron generation time.

    \vspace{1em}

    \begin{alertblock}{The Speed Problem}

        $\Lambda$ is far shorter than human or ordinary plant-control
        timescales.

        Delayed neutrons are therefore essential to normal reactor
        controllability.

    \end{alertblock}

\end{frame}


% -----------------------------------------------------------------------------
% Delayed neutrons
% -----------------------------------------------------------------------------

\begin{frame}{Step 3: Delayed Neutrons}

    A small fraction $\beta$ of fission neutrons are emitted after radioactive
    precursor decay \cite{bell1970nuclear}.

    \vspace{1em}

    \begin{itemize}

        \item
        \textbf{Prompt neutrons:}
        born essentially with the fission event.

        \item
        \textbf{Delayed neutrons:}
        emitted from precursor decay over much longer timescales.

    \end{itemize}

    \vspace{1em}

    The precursor inventory gives the reactor kinetic
    \textbf{memory} observed in the simulator.

\end{frame}


% -----------------------------------------------------------------------------
% PRKE
% -----------------------------------------------------------------------------

\begin{frame}{The Point Reactor Kinetics Equations}

    For $G$ delayed-neutron groups:

    \begin{align*}
        \frac{dn}{dt}
        &=
        \frac{\rho-\beta}{\Lambda}n
        +
        \sum_{i=1}^{G}\lambda_i C_i
        \\[1em]
        \frac{dC_i}{dt}
        &=
        \frac{\beta_i}{\Lambda}n
        -
        \lambda_i C_i,
        \qquad
        i=1,\ldots,G
        \\[1em]
        \beta
        &=
        \sum_{i=1}^{G}\beta_i
    \end{align*}

\end{frame}


% -----------------------------------------------------------------------------
% Five group implementation
% -----------------------------------------------------------------------------

\begin{frame}{Delayed Neutrons in HTGR Sim v1}

    The textbook PRKE is commonly introduced using six delayed-neutron groups.

    \vspace{1em}

    \begin{block}{HTGR Sim v1}

        The current demonstrator uses a reduced
        \textbf{five-group U-235 delayed-neutron bank},
        coupled to its prompt excursion model.

    \end{block}

    \vspace{1em}

    The group structure is a modelling choice.

    The underlying physical idea remains the same:

    \[
        \text{fission}
        \rightarrow
        \text{precursors}
        \rightarrow
        \text{delayed neutrons}
    \]

\end{frame}


% -----------------------------------------------------------------------------
% Multiple groups
% -----------------------------------------------------------------------------

\begin{frame}{Why Multiple Delayed Groups Matter}

    Different precursor families decay on different timescales.

    \vspace{1em}

    \begin{itemize}

        \item Fast precursor groups influence the early delayed response.

        \item Slow groups preserve reactor-history information for much longer.

        \item A single effective group can teach the basic concept.

        \item Multiple groups provide a more faithful transient response.

    \end{itemize}

    \vspace{1em}

    HTGR Sim v1 therefore retains a reduced multi-group precursor bank rather
    than using prompt kinetics alone.

\end{frame}


% =============================================================================
% PART III -- REACTIVITY
% =============================================================================

\section{Reactivity and Feedback}


% -----------------------------------------------------------------------------
% Prompt critical
% -----------------------------------------------------------------------------

\begin{frame}{The Edge of Control: Prompt Criticality}

    The prompt contribution contains:

    \[
        \frac{\rho-\beta}{\Lambda} n
    \]

    \vspace{1em}

    \begin{itemize}

        \item For $\rho < \beta$,
              delayed neutrons remain essential to the increasing solution.

        \item At $\rho = \beta$,
              the reactor is \textbf{prompt critical}.

        \item For $\rho > \beta$,
              prompt-neutron dynamics alone can sustain rapid growth.

    \end{itemize}

    \begin{alertblock}{Important}

        Prompt criticality marks a qualitative change in the governing
        kinetics.

        It is not a normal control regime.

    \end{alertblock}

\end{frame}


% -----------------------------------------------------------------------------
% Dollars
% -----------------------------------------------------------------------------

\begin{frame}{Units of Reactivity: Dollars}

    We scale reactivity by the effective delayed-neutron fraction:

    \[
        \$
        =
        \frac{\rho}{\beta}
    \]

    \vspace{1em}

    \begin{itemize}

        \item $\$0$:
              critical reference condition

        \item positive cents:
              positive reactivity below prompt critical

        \item $\$1$:
              $\rho = \beta$, the prompt-critical boundary

    \end{itemize}

    \begin{block}{HTGR Sim v1}

        The GUI drives reactivity in dollars.

        The kinetics layer converts this to dimensionless external reactivity
        before advancing the model.

    \end{block}

\end{frame}


% -----------------------------------------------------------------------------
% Challenge 2
% -----------------------------------------------------------------------------

\begin{frame}{Challenge Redux: Watch the Dollars}

    Repeat the power manoeuvre.

    This time, watch the reactivity display.

    \vspace{1em}

    \begin{itemize}

        \item How much positive reactivity is actually needed?

        \item How long must it remain positive?

        \item What happens when external reactivity returns toward zero?

        \item How does the temperature-feedback contribution change during the
              transient?

    \end{itemize}

    \vspace{1em}

    \centering

    \textbf{
        Can you now anticipate the reactor response before it happens?
    }

\end{frame}


% -----------------------------------------------------------------------------
% Feedback
% -----------------------------------------------------------------------------

\begin{frame}{Why Doesn't Power Increase Forever?}

    External control is only part of the total reactivity:

    \[
        \rho(t)
        =
        \rho_{\mathrm{external}}
        +
        \rho_{\mathrm{feedback}}(T,\ldots)
    \]

    \vspace{1em}

    HTGR Sim v1 includes an illustrative
    \textbf{negative fuel-temperature coefficient}.

    As the lumped fuel/core temperature rises, this contribution becomes more
    negative and opposes the excursion.

    \vspace{1em}

    \begin{block}{Physical Idea}

        Reactor kinetics and thermal hydraulics are coupled:

        \[
            \text{Power}
            \rightarrow
            \text{Temperature}
            \rightarrow
            \text{Reactivity}
            \rightarrow
            \text{Power}
        \]

    \end{block}

\end{frame}


% -----------------------------------------------------------------------------
% Thermal inertia
% -----------------------------------------------------------------------------

\begin{frame}{HTGR Thermal Inertia}

    An HTGR contains a large graphite inventory and operates with helium
    coolant.

    \vspace{1em}

    \begin{itemize}

        \item Neutron power can change rapidly.

        \item The graphite/fuel thermal state changes much more slowly.

        \item Helium transports heat from the pebble bed toward the
              steam generator.

        \item Thermal feedback therefore evolves on a different timescale from
              the prompt neutron population.

    \end{itemize}

    \vspace{1em}

    This separation of timescales is one reason an interactive coupled
    simulator is useful for teaching reactor dynamics.

\end{frame}


% =============================================================================
% PART IV -- DECAY HEAT
% =============================================================================

\section{Shutdown and Decay Heat}


% -----------------------------------------------------------------------------
% Shutdown
% -----------------------------------------------------------------------------

\begin{frame}{Can We Switch a Reactor Off?}

    Inserting sufficiently negative reactivity suppresses the fission chain
    reaction.

    But the reactor does \textbf{not} instantly become thermally cold.

    \vspace{1em}

    Fission products continue to decay and release heat.

    \vspace{1em}

    \begin{alertblock}{Key Distinction}

        \textbf{Fission power}
        can collapse rapidly after shutdown.

        \vspace{0.5em}

        \textbf{Decay heat}
        persists and must still be removed.

    \end{alertblock}

\end{frame}


% -----------------------------------------------------------------------------
% Decay heat model
% -----------------------------------------------------------------------------

\begin{frame}{Decay Heat in HTGR Sim v1}

    The demonstrator includes a multi-group fission-product decay-heat model.

    \vspace{1em}

    Its thermal source is conceptually:

    \[
        P_{\mathrm{thermal}}
        =
        P_{\mathrm{prompt/fission\ component}}
        +
        P_{\mathrm{decay}}
    \]

    \vspace{1em}

    The decay-heat bank is initialised consistently with an operating reactor,
    rather than assuming a fresh core with no accumulated fission products.

    \begin{block}{Teaching Consequence}

        A shutdown demonstration can show why

        \begin{center}
            ``reactor shut down''
            $\neq$
            ``heat generation equals zero''.
        \end{center}

    \end{block}

\end{frame}


% -----------------------------------------------------------------------------
% Shutdown challenge
% -----------------------------------------------------------------------------

\begin{frame}{Shutdown Demonstration}

    \begin{columns}

        \begin{column}{0.52\textwidth}

            \textbf{Try it:}

            \begin{enumerate}

                \item Establish a stable operating condition.

                \item Insert strong negative reactivity or perform the
                      simulator shutdown action.

                \item Watch fission power fall.

                \item Continue watching core temperature and thermal power.

                \item Identify the slower decay-heat tail.

            \end{enumerate}

        \end{column}

        \begin{column}{0.48\textwidth}

            \centering

            \IfFileExists{\htgrshutdown}
            {
                \includegraphics[width=\textwidth]{\htgrshutdown}
            }
            {
                \fbox{
                    \parbox[c][0.45\textheight][c]{0.9\textwidth}{
                        \centering
                        Insert HTGR shutdown screenshot or decay-heat plot here.
                    }
                }
            }

        \end{column}

    \end{columns}

\end{frame}


% -----------------------------------------------------------------------------
% Shutdown time history
% -----------------------------------------------------------------------------

\begin{frame}{Shutdown: What the Time History Shows}

    \centering

    \IfFileExists{\htgrshutdownplot}
    {
        \includegraphics[
            width=\textwidth,
            height=0.70\textheight,
            keepaspectratio
        ]{\htgrshutdownplot}
    }
    {
        \fbox{
            \parbox[c][0.55\textheight][c]{0.82\textwidth}{
                \centering
                Insert HTGR Sim v1 time-history plot here.
            }
        }
    }

    \vspace{0.6em}

    \raggedright
    \footnotesize
    Rod insertion drives fission power down sharply, but thermal power does not
    follow it to zero. The tail is decay heat, and it is why a shut-down reactor
    still needs cooling.

\end{frame}


% =============================================================================
% PART V -- MODEL FIDELITY
% =============================================================================

\section{What the Demonstrator Represents}


% -----------------------------------------------------------------------------
% What it models
% -----------------------------------------------------------------------------

\begin{frame}{What HTGR Sim v1 Actually Models}

    \begin{itemize}

        \item Point-reactor power dynamics.

        \item Prompt excursion / temperature-feedback layer.

        \item Reduced five-group U-235 delayed-neutron precursor dynamics.

        \item Fission-product decay heat.

        \item Lumped pebble-bed/core thermal response.

        \item Helium-loop and steam-generator demonstration physics.

        \item Coupling between reactor kinetics and coolant heat removal.

    \end{itemize}

\end{frame}


% -----------------------------------------------------------------------------
% Limitations
% -----------------------------------------------------------------------------

\begin{frame}{Model Limitations}

    \begin{alertblock}{Point Kinetics}

        The core has one neutron-power amplitude.

        There is no spatial flux shape, axial power distribution or local
        power peak.

    \end{alertblock}

    \begin{itemize}

        \item Control-rod worth does not depend on individual rod position.

        \item Fuel-temperature feedback uses a lumped temperature, not resolved
              TRISO fuel-kernel temperatures.

        \item Plant-scale coefficients and controller parameters include
              illustrative demonstration assumptions.

        \item The model is not validated for reactor operation, licensing,
              safety analysis or operator training.

    \end{itemize}

\end{frame}


% -----------------------------------------------------------------------------
% Fidelity hierarchy
% -----------------------------------------------------------------------------

\begin{frame}{From Point Kinetics to Higher Fidelity}

    The point model deliberately trades spatial fidelity for speed and clarity.

    \vspace{1em}

    \[
        \boxed{\text{Point Kinetics}}
        \quad\longrightarrow\quad
        \boxed{\text{Coarse Nodal Neutronics}}
        \quad\longrightarrow\quad
        \boxed{\text{High-Fidelity Transport}}
    \]

    \vspace{1.5em}

    In a digital-twin architecture, different fidelity levels can serve
    different timescales and questions.

    \vspace{1em}

    A fast model is valuable precisely because it can run interactively
    \textbf{provided its limitations are explicit}.

\end{frame}


% =============================================================================
% PART VI -- TAKEAWAYS
% =============================================================================

\section{Takeaways}


% -----------------------------------------------------------------------------
% Summary
% -----------------------------------------------------------------------------

\begin{frame}{What You Should Take Away}

    \begin{enumerate}

        \item Reactor power follows neutron-population dynamics.

        \item Delayed neutrons make normal reactor control possible on
              engineering timescales.

        \item Reactivity in dollars makes the prompt-critical boundary
              intuitive.

        \item Temperature feedback couples neutronics to thermal hydraulics.

        \item HTGR thermal inertia is much slower than prompt neutron kinetics.

        \item Shutdown ends the chain reaction, but
              \textbf{decay heat remains}.

        \item A simulator is only as trustworthy as the physics and validation
              supporting each model.

    \end{enumerate}

\end{frame}


% -----------------------------------------------------------------------------
% Final challenge
% -----------------------------------------------------------------------------

\begin{frame}{One More Run}

    \centering

    \Large

    Now operate the simulator again.

    \vspace{1.5em}

    \textbf{
        This time, explain the transient while it happens.
    }

    \vspace{1.5em}

    \normalsize

    What are the prompt neutrons doing?

    \vspace{0.5em}

    What are the precursors doing?

    \vspace{0.5em}

    What is temperature feedback doing?

    \vspace{0.5em}

    Where does the heat go after shutdown?

\end{frame}


% =============================================================================
% REFERENCES
% =============================================================================

% `allowframebreaks` segfaults LuaTeX with this theme's frametitle template,
% which corrupts the .bcf and so breaks biber. Do not reintroduce it here.
\begin{frame}{References}

    \printbibliography

\end{frame}


% =============================================================================
% END
% =============================================================================

\begin{frame}

    \centering

    \Huge
    \textbf{Questions?}

    \vspace{1em}

    \Large

    Point Kinetics
    $\longrightarrow$
    Coupled Reactor Dynamics

\end{frame}


\end{document}
